%======================================================================
\chapter{Stability: first notions}\label{ch:st}
%======================================================================
% Template chapter: the frame of every chapter of the book
% (Outline, numbered sections, Summary, Exercises, References).

\framepart{Outline}
\draftnote{One paragraph: what the chapter does and why.}

\paragraph{By the end of this chapter you should be able to:}
\begin{itemize}
  \item \draftnote{objective 1;}
  \item \draftnote{objective 2.}
\end{itemize}

\paragraph{Plan of the chapter.} Section~\ref{sec:st-intro} \draftnote{plan}.
A summary, exercises with solutions (page~\pageref{sec:st-exo}) and the
references close the chapter.

\section{Introduction}\label{sec:st-intro}
\index{stability}%
\draftnote{Start with a small example with numbers before any general statement.}

\begin{gbox}{A key statement, written as a claim in words}\label{sec:st-key}
\draftnote{Key result in a gray box; cite with Section~\ref{sec:st-key} or \nameref{sec:st-key}.}
\end{gbox}
\paragraph{Why it works.}
\draftnote{The picture first, then the short chain of formulas.}

%======================================================================
\framepart{Summary}
\begin{gbox*}
\begin{itemize}
  \item \textbf{First point.} \draftnote{take-home message.}
\end{itemize}
\end{gbox*}

%======================================================================
\framepart{Exercises}\label{sec:st-exo}

\framesub{Short questions}

\exercise{A first question}\label{exo:st-first}
\draftnote{Statement.}
\begin{solution}
\draftnote{Short solution.}
\end{solution}

\framesub{Problems}

%======================================================================
\begin{chapterbib}{99}
\bibitem{st-placeholder} \draftnote{First reference of Chapter 1 (keys unique book-wide).}
\end{chapterbib}
